\documentclass[11pt]{article}
\usepackage[utf8]{inputenc}
\usepackage{csquotes}
\usepackage{setspace}
\usepackage[dvipsnames]{xcolor}
\usepackage[margin=1in]{geometry}
\usepackage{hyperref}
\usepackage{amsmath, amssymb, amscd, amsthm, amsfonts}
\usepackage{graphicx}
\usepackage{adjustbox}
\usepackage{hyperref}
\usepackage{cleveref}
\usepackage{amsthm}
\usepackage[utf8]{inputenc}
\usepackage[english]{babel}
\usepackage{proof}
\usepackage{ninecolors}
\usepackage{enumitem}

\newcommand{\B}{\mathbf{B}}
\newcommand{\w}{\mathbf{w}}
\usepackage{proof}
\usepackage{tikz-cd}
\usepackage{xcolor}
\usepackage{thmtools, thm-restate}
\tikzcdset{scale cd/.style={every label/.append style={scale=#1},
		cells={nodes={scale=#1}}}}
%\usepackage{stmaryrd}


\urlstyle{same}

\newtheorem{theorem}{Theorem}[section]
\newtheorem{corollary}{Corollary}[theorem]
\newtheorem{lemma}[theorem]{Lemma}
\newtheorem{conjecture}[theorem]{Conjecture}
\newtheorem{proposition}[theorem]{Proposition}

\declaretheorem[name=Theorem,numberwithin=section]{thm}

\theoremstyle{definition}
\newtheorem{definition}{Definition}[section]

\title{Simulations of Simplicial Actions for Distributed Protocols}
\author{Natalia Meija Bautista, Philip Sink, Taiba Abid}
\date{}

\newcommand{\rr}{\mathbb{R}}

\newcommand{\al}{\alpha}
\DeclareMathOperator{\conv}{conv}
\DeclareMathOperator{\aff}{aff}
\newcommand{\V}{\mathcal{V}}
\newcommand{\M}{\mathcal{M}}
\newcommand{\td}{\bigtriangledown}
\newcommand{\tu}{\bigtriangleup}

%ADAM'S MACROS

\newcommand{\draft}[1]{{\color{red}{\textsc{[#1]}}}}
\newcommand{\commentout}[1]{}
\newcommand{\defin}[1]{\textbf{#1}}

\renewcommand{\phi}{\varphi}

\newcommand{\lthen}{\rightarrow}
\newcommand{\liff}{\leftrightarrow}

\newcommand{\F}{\mathcal{F}}
\renewcommand{\M}{\mathcal{M}}
\renewcommand{\L}{\mathcal{L}}
\newcommand{\N}{\mathcal{N}}

\newcommand{\full}{\textsf{FULL}}

%a better set-minus%
\usepackage{amssymb,tikz}

\newcommand{\mysetminusD}{\hbox{\tikz{\draw[line width=0.6pt,line cap=round] (3pt,0) -- (0,6pt);}}}
\newcommand{\mysetminusT}{\mysetminusD}
\newcommand{\mysetminusS}{\hbox{\tikz{\draw[line width=0.45pt,line cap=round] (2pt,0) -- (0,4pt);}}}
\newcommand{\mysetminusSS}{\hbox{\tikz{\draw[line width=0.4pt,line cap=round] (1.5pt,0) -- (0,3pt);}}}

\newcommand{\mysetminus}{\mathbin{\mathchoice{\mysetminusD}{\mysetminusT}{\mysetminusS}{\mysetminusSS}}}

\begin{document}
	
	\maketitle
	
	\begin{abstract}
		

		
	\end{abstract}
	
	\section{Introduction} \label{sec:int}
	
	Seeing as the impetus for Simplicial Semantics is the use of combinatorial topology in distributed computing, most of that literature acknowledges and is indebted to this connection already. \cite{Death,DoA,KaSC,SimpDEL,HG,FA,DCTCT,CACHIN2025114902,Roho_2025,Lehnherr_2025} In general, the procedures for using action models to describe distributed protocols are understood but new, showing up in \cite{SimpDEL}, \cite{KaSC}, \cite{KaG}, \cite{ActionDC}, \cite{CPM}, and \cite{CPaAM}. Moreover, the papers \cite{KaSC}, \cite{BelSimp}, \cite{HGD}, and \cite{SimpBel} are already attempts to give a semantics for belief, or non-factive (defeasible) knowledge. In \cite{TechMemo} and \cite{SimpAct}, these threads were brought together by using action models in the setting of simplicial semantics to model distributed protocols, and then incorporating belief revision, as defined in \cite{Lewis2} and \cite{Stalnaker1981}, and implemented in the simplicial setting as in \cite{SimpBelRev}, to then define new protocols and a new model of communication, in the sense of \cite{TopPersp}. For those unfamiliar with the details of this work, we will summarize it in Section \ref{sec:back}.
	
	\section{Background}\label{sec:back}
	
	The goal of this project is to use Python to simulate the definitions given in \cite{TechMemo} and \cite{SimpAct}. This section will give the particulars for how we will define simplicial semantics and simplicial actions throughout the paper. For a more thorough overview, we recommend reading \cite{KaSC}, \cite{BelSimp}, or \cite{SimpBelRev}. For more on actions specifically, one can consult \cite{TechMemo} or \cite{SimpAct}.
	
	The literature on simplicial semantics divides roughly into two approaches to defining valuations for propositional atoms. Loosely speaking, the first approach assigns truth values directly to the facets \cite{Death,FA,SimpSet},	while the second ``vertex-based'' approach assigns them (in a partial way) to the vertices and then ``lifts'' them to facets. \cite{DoA,KaSC,SimpDEL,HG} This paper will focus on the latter approach.
	
	Let $\mathfrak{P}$ be a countable set of propositional atoms, and $Ag$ a finite set of agents. Let $N$ be a set of of \defin{nodes}, $V:\rightarrow Ag$ a function called the \defin{coloring} function, and $L:N\rightarrow 3^\mathfrak{P}$ a function which assigns each node to a set of literals, which we call the \defin{assignment}. The interpretation is that $L(n)(P)=1$ if and only if $P$ is associated with $n$, $L(n)(P)=0$ if and only if $\neg P$ is associated with $n$, and $L(n)(P)=2$ if and only if neither is associated with $n$. $L$ is the only difference between their presentation here as compared to \cite{BelSimp}. As mentioned above, $L$ assigns propositions to nodes.
	
	Our first key idea is that we can use $N$, $V$, and $L$ to create a kind of \defin{maximal} simplicial complex. Like much of the previous literature, we will assume our simplicial complexes are uniquely colored. Specifically, each facet of our complexes has a dimension of size $|Ag|$, and no two nodes are associated to the same agent. Put formally, if $X\in S$ is such that for all $Y\in S$, if $X\subseteq Y$ then $X=Y$, we have that $|X|=|Ag|$, and, if $x,y\in X$, then if $V(x)=V(y)$, we have that $x=y$. We call this condition UCF for ``Uniquely Colored Facets''. The \defin{Maximal Complex} is the set of subsets of $N$ such that the associated logical content with that subset is consistent, and it satisfies the UCF condition. That is, it's the subsets $x\subseteq N$ such that $|x|=|A|$, The set of formulas assigned to $x$ by $L$ is consistent, and for all $u,v\in x$, $V(u)\neq V(v)$. We refer to $\mathfrak{M}((N,V,L)$ as the \defin{Maximal Complex} of $N$, $V$, and $L$. When the context is clear, we will refer to it simply as the maximal complex. The maximal complex can be defined set theoretically as follows:
	
	\begin{align*}
		\mathfrak{M}(N,V,L):=\{y\in 2^N~|~\exists x\in 2^N (&y\subseteq x\\&\wedge\neg(\exists P\in\mathfrak{P}(\exists u,v\in x(L(u)(P)=1\wedge L(v)(P)=0))) \\&\wedge (|x|=|Ag|)
		\\&\wedge (\forall u,v\in x(V(u)\neq V(v))))\}
	\end{align*}
	
	Note that $\mathfrak{M}(N,V,L)$ is a simplicial complex. This is because, if $Y\in \mathfrak{M}(N,V,L)$, there is a set $X$ such that $Y\subseteq X$ and $X$ satisfies certain properties. If $Z\subseteq Y$, then the same set $X$ suffices to show that $Z\in \mathfrak{M}(N,V,L)$. All complexes we consider in this paper will be UCF subcomplexes of the maximal complex. 
	
	If $S$ is a simplicial complex, let $\mathcal{F}(S)$ denote the facets of $S$. Since all simplicial complexes we consider are UCF, we can define projection functions $\pi_a:\mathcal{F}(S)\rightarrow N$ given by $\pi(a)(X)=V^{-1}(a)\cap X$. That is, $\pi_a$ maps each facet to the unique $a$-colored node it contains.
	
	We can now define a \defin{Simplicial Model For Knowledge}:
	
	\begin{definition}[Simplicial Model for Knowledge]
		A \defin{Simplicial Model for Knowledge} $\M$ is a tuple $(N,V,L,S)$ where $N$ is a nonempty set of nodes, $V$ is a coloring function, $L$ is an assignment function, and $S$ is a UCF subcomplex of $\mathfrak{M}(N,V,L)$.
	\end{definition}
	
	In this setting, the language $\L_{K}(Ag)$ recursively defined by
	$$\varphi ::= P \, | \, \bot \, | \, \phi \lthen \psi \, | \, K_{a}\phi,$$
	where $P \in \mathfrak{P}$ and $a \in Ag$, can be interpreted in simplicial models for knowledge as follows for any $X\in\F(S)$:
	\begin{align*}
		\M,X & \models P \text{ iff } \exists a\in Ag(L(\pi_a(X))(P)=1)\\
		\M,X & \nvDash \bot\\
		\M,X & \models \phi \lthen \psi \text{ iff } \M,X \models \phi \text{ implies } \M,X \models \psi\\
		\M,X & \models K_a \phi \text{ iff } \forall Y \in \F(S) \text{ if } \pi_a(Y) = \pi_a(X) \text{ then } \M,Y \models \phi
	\end{align*}
	
	In order to incorporate belief revision, it will be necessary to talk about models for belief rather than models for knowledge. The following definition is borrowed from \cite{SimpBelRev}:
	
	\begin{definition}[Simplicial Model for Belief]
		A \defin{Simplicial Model for Belief} $\M$ is a tuple $(N,V,L,S,\{S_a\}_{a\in Ag})$ where $N$, $V$, $L$, and $S$ are familiar, and each $S_a$ is a UCF subcomplex of $S$. Often, we will leave $S$ unspecified. In this case, we assume that $S=\mathfrak{M}(N,V,L)$.
	\end{definition}
	
	In this setting, the language $\L_{B}(Ag)$ recursively defined by
	$$\varphi ::= P \, | \, \bot \, | \, \phi \lthen \psi \, | \, B_{a}\phi,$$
	where $P \in \mathfrak{P}$ and $a \in Ag$, can be interpreted in simplicial models for knowledge as follows for any $X\in\F(S)$:
	\begin{align*}
		\M,X & \models P \text{ iff } \exists a\in Ag(L(\pi_a(X))(P)=1)\\
		\M,X & \nvDash \bot\\
		\M,X & \models \phi \lthen \psi \text{ iff } \M,X \models \phi \text{ implies } \M,X \models \psi\\
		\M,X & \models B_a \phi \text{ iff } \forall Y \in \F(S_a) \text{ if } \pi_a(Y) = \pi_a(X) \text{ then } \M,Y \models \phi
	\end{align*}
	
	Following common convention, if $\M$ is a simplicial model for either knowledge or belief, the claim $x\in\M$ is short hand for $x\in N$ where $N$ is the set of nodes in $\M$.
	
	The final notion is a \defin{Simplicial Model for Belief and Knowledge}. This is a tuple $(N,V,L,S,\{S_a\}_{a\in Ag})$, and, as in the belief case, when $S$ is unspecified, we assume that $S=\mathfrak{M}(N,V,L)$. 	In this setting, the language $\L_{KB}(Ag)$ recursively defined by
	$$\varphi ::= P \, | \, \bot \, | \, \phi \lthen \psi \, | \, K_{a}\phi \, | \, B_{a}\phi$$
	where $P \in \mathfrak{P}$ and $a \in Ag$, can be interpreted in simplicial models for knowledge as follows for any $X\in\F(S)$. One appeals to $S$ for the truth condition for $K_a$ as in simplicial models for knowledge, and to the $S_a$ for the truth condition for $B_a$ as in simplicial models for belief.
	
	We say that a relational model is called \defin{proper} if, for all $w \in W$, $\big|\bigcap_{a \in Ag} [w]_{a} \big| = 1$, and for all $P$ in our language and any $w\in W$, there is an $a\in Ag$ such that for all $u\in R_a(w)$, $u\in v(P)$. Call the set of such distinguished agents $d(w,P)$.\footnote{This definition of properness differs from standard presentations when the relations are not equivalence relations, as has also been noted in \cite{Lehnherr_2025}. It also differs from the presentation in \cite{BelSimp} by adding in the distinguished agent. This is because there, unlike here, atmoic propositions are assigned to facets, not nodes.} Note that when a relational model is proper, the following axioms are sound for any atomic formula $P$:
	
	$$\textsf{NUK}:P\rightarrow\bigvee_{a\in Ag}K_aP$$
	
	$$\textsf{NUB}:P\rightarrow\bigvee_{a\in Ag}B_aP$$
	
	Here, ``NU'' stands for ``No Uncertainties''. The idea is that an atomic formula cannot be true unless there is some agent which knows/believes it. This is clearly sound in simplicial models, so the restriction of relational models to those which are proper, and hence satisfy \textsf{NUK} and \textsf{NUB}, is necessary in order to translate between the two.
	
	This translation is another key result we simulate. Although the implementation of action models and especially belief revision in \cite{SimpBelRev}, \cite{SimpAct}, and \cite{TechMemo} depend heavily on the simplicial semantics, the much wider use of relational models\footnote{See \cite{DC5} for an overview, and \cite{KaSC}, \cite{KaG}, \cite{ActionDC}, \cite{CPM}, and \cite{CPaAM} for approaches closely aligned with ours, in particular in the use of action models.} means that some users may find it easier to define a model relationally and then translate it into a simplicial model to take advantage of these tools. As such, we introduce relational models, and in particular how to translate them into simplicial models, as in \cite{BelSimp}.
	
	\begin{definition}[Relational Model for Knowledge] A \defin{Relational Model (over $Ag$)} is a tuple $\N = (W,(R_{a})_{a \in Ag},v)$ where $W$ is a nonempty set of \emph{possible worlds}, each $R_{a}$ is a binary relation on $W$ called an \emph{accessibility relation}, and $v: \mathfrak{P} \to 2^{W}$ is a \emph{valuation}. Semantics for the language $\L_{K}(Ag)$ are given in the usual way:
	\begin{align*}
		\N,w & \models P \text{ iff } w \in v(P)\\
		\N,w & \nvDash \bot\\
		\N,w & \models \phi \lthen \psi \text{ iff } \N,w \models \phi \text{ implies } \N,w \models \psi\\
		\N,w & \models K_a \phi \text{ iff } \forall w' \in R_{a}(w) (\N,w' \models \phi),
	\end{align*}
	where $R_{a}(w) = \{w' \in W \: : \: wR_{a}w'\}$. When each $R_{a}$ is an equivalence relation or partition, we call this a \defin{Relational Model for Knowledge}. Such models will satisfy \defin{Factivity} ($\vdash K_a\varphi\rightarrow\varphi$), \defin{Positive Introspection} ($\vdash K_a\varphi\rightarrow K_aK_a\varphi$), and \defin{Negative Introspection} ($\vdash \neg K_a\varphi\rightarrow K_a\neg K_a\varphi$). Respectively, we call these axioms \textsf{T}, \textsf{4}, and \textsf{5}.
	\end{definition}

	We also have relational models for belief.
	
		\begin{definition}[Relational Model for Belief] Let $\N = (W,(Q_{a})_{a \in Ag},v)$ be a relational model. Semantics are given in the usual way, though for the language $\L_{B}(Ag)$. When each $Q_a$ is transitive Euclidean\footnote{For any worlds $w$, $u$, and $v$, if $u,v\in Q_a(w)$ then $u\in Q_a(v)$. Note that relational models for knowledge are also Euclidean.}, we call this a \defin{Relational Model for Belief}. Such models will satisfy positive introspection and negative introspection.\footnote{Typical work on belief, such as \cite{HGD}, \cite{SimpBel}, and \cite{BelSimp} uses the \textsf{KD45} axiom system, which adds \textsf{D}, consistency, ensuring $\neg B_a\bot$. However, as observed in \cite{SimpBelRev}, simplicial action models occasionally isolate vertices in such a way that \textsf{D} cannot be preserved in outputs, even if it is true in inputs and actions. So, we do not include it here.}
	\end{definition}

	The final case is when both belief and knowledge are present. A \defin{Relational Model for Belief and Knowledge} is a tuple $\N=(W,(R_{a})_{a \in Ag},(Q_{a})_{a \in Ag},v)$, where each $R_a$ is an equivalence relation, each $Q_a$ is serial, transitive, and Euclidean, $Q_a\subseteq R_a$ for each $a\in Ag$, and each $Q_a$ is constant on $R_a$.\footnote{If $u\in R_a(w)$, then $Q_a(w)=Q_a(u)$.} Such models can interpret the language $\L_{KB}(Ag)$ by using the $R_a$ to interpret $K_a\varphi$, just as in relational models for knowledge, and using the $Q_a$ to interpret $B_a\varphi$, just as in relational models for belief. The following axioms will be sound in such models:
	
	\begin{itemize}
	\item knowledge implies belief ($K_{a} \phi \lthen B_{a} \phi$)
	\item strong positive introspection ($B_{a} \phi \lthen K_{a} B_{a} \phi$)
	\item strong negative introspection ($\lnot B_{a} \phi \lthen K_{a} \lnot B_{a} \phi)$
	\end{itemize}

	In fact, these are sufficient. The following is well known: \cite{DC5, BelSimp,Bjorndahl2024}
	
	\begin{theorem}[Soundness and Completeness for Relational Models]
	The smallest axiom system closed under propositional logic, \textsf{T}, \textsf{4}, \textsf{5}, \textsf{K}\footnote{$\vdash ((K_a\varphi\wedge K_a(\varphi\rightarrow\psi)\rightarrow K_a\psi)$} and the rules of inference modus ponens, necessitation\footnote{$\vdash \varphi\Rightarrow\vdash K_a\varphi$} is called $\textsf{S5}$. Relational Models for knowledge are sound and complete with respect to the axiom system $\textsf{S5}$. The smallest axiom system closed under propositional logic, \textsf{4}, \textsf{5}, \textsf{K} and the rules of inference modus ponens, necessitation is called $\textsf{K45}$. Relational Models for belief are sound and complete with respect to the axiom system $\textsf{K45}$. Let $\full$ be the axiom system for the language $\L_{KB}(Ag)$ obtained by adding the above three axiom schemes to the system \textsf{S5} for knowledge combined with \textsf{K45} for belief. Relational Models for Knowledge and Belief are sound and complete with respect to the axiom system \full. When any of these classes are proper, they are sound and complete with respect to the same axiom systems simply with the schemes \textsf{NUK} and \textsf{NUB} added in as appropriate for the relevant language.
	\end{theorem}

	We can now turn to translating relational models into simplicial models. The following proposition is from \cite{BelSimp}:
	
	\begin{proposition}
		Let $\N$ be a relational model for belief and knowledge. Using the L\"owenheim-Skolem theorem, fix $G$ a group of order $|W|$\footnote{Following convention, we denote the operation with concatenation and inverse with $()^{-1}$} with, and $g:W\rightarrow G$ a bijection. Moreover, let $i\in Ag$ be a distinguished agent. Define $\tilde{N}$ to be the relational model $(\tilde{W}, (\tilde{R}_a)_{a\in Ag}, (\tilde{Q}_a)_{a\in Ag}, v)$, where
		\begin{align*}
		\tilde{W}&:=W^2\\
		(x',y')\in\tilde{R}_a(x,y)&\text{ iff }x'\in R_a(x)\text{ and } \begin{cases}
		g(y)g(x)^{-1} = g(y')g(x')^{-1}& \text{ if }\;a=i\\
		y=y'& \text{ otherwise}
		\end{cases}\\
		(x',y')\in\tilde{Q}_a(x,y)&\text{ iff }x'\in Q_a(x)\text{ and } \begin{cases}
			g(y)g(x)^{-1} = g(y')g(x')^{-1}& \text{ if }\;a=i\\
			y=y'& \text{ otherwise}
		\end{cases}\\
	\tilde{v}(P)&:=\{(x,y) \in \tilde{W} \: : \: x \in v(P)\}
		\end{align*}
		Then $\tilde{\N}$ is a proper relational model for knowledge and belief, and projection to the first coordinate is a surjective, bounded morphism from $\tilde{\N}$ to $\N$.
	\end{proposition}

	The obvious parallel results for relational models for only knowledge, or only belief, also hold.
	
	Only proper relational models can be transformed into simplicial models. Hence, the above proposition allows us to translate any relational model for knowledge into a simplicial model.	Given a proper relational model for knowledge and belief $\N=(W,(R_a)_{a\in Ag},(Q_a)_{a\in Ag},v)$, we define a simplicial model for knowledge and belief $\M_{\N}:=(N_{\N},V_{\N},L_{\N},S_{\N},(S_{\N,a})_{a\in Ag})$ similar to as in \cite{BelSimp}:\footnote{The difference between the work presented here, and the work in \cite{BelSimp}, is that in the latter, atomic formulae are assigned to facets, and here, they are assigned to nodes.}\footnote{Observe that we do not need to use the 0 case for $L_\N$. This is because the truth value of atomic propositions does not appeal to this case. It only affects what the maximal complex is, but because we are specifying $S$ directly, this is not relevant.}
	
\begin{align*}
		N_{\N} &:= \{([w]_{a},a) \: : \: w \in W, a \in Ag\}\\
V_{\N}([w]_{a},a) &:= a\\
L_{\N}([w]_{a},a)(P) &:=\begin{cases}
	1& \text{ if }\;a\in d(w,P)\\
	2& \text{ otherwise}
\end{cases}\\
S_{\N} &:= \{F \subseteq N_{\N} \: : \: \bigcap_{n \in F} e(n) \neq \emptyset\}
\end{align*}

	where $e: N_{\N} \to 2^{W}$ is projection to the first coordinate (returning the equivalence class associated with each node), and $f \colon W \to \F(S_{\N})$ is the bijective correspondence between worlds in $W$ and facets in $S_{\N}$. What's missing are the belief subcomplexes, which we define as follows:
	$$\mathcal{F}(S_{\N,a}) := \{F\in\mathcal{F}(S_{\N}) \: : \: f^{-1}(F) Q_{a} f^{-1}(F)\},$$
	Specifying the facets of $S_{\N,a}$ uniquely determines the simplicial complex itself; it also guarantees that $S_{\N,a}$ is a UCF subcomplex of $S_{\N}$. 
	
	Translations similar to this are not new to the literature, appearing in \cite{SimpDEL} and \cite{KaSC}. Analogous translations for the case of relational/simplicial models with only knowledge or only belief are obvious. Such a translation gives us the following propositions, as in \cite{BelSimp}:
	
	\begin{proposition}
		The simplicial frame $(N_{\N},V_{\N},L_{\N},S_{\N},(S_{\N,a})_{a\in Ag})$ satisfies UCF.
	\end{proposition}

	\begin{proposition}
		For any relational model for knowledge and belief $\N$, the simplicial frame $\M_{\N}$ is such that $\N,w\vDash\varphi$ if and only if $\M_\N,f(w)\vDash\varphi$.
	\end{proposition}

	Which, in turn, gives us the following theorem(s) as in \cite{BelSimp} and \cite{SimpBelRev}:
	
	\begin{theorem}[Soundness and Completeness for Simplicial Models]
		Simplicial Models for knowledge and belief (knowledge) (belief) are sound and complete with respect to \full+\textsf{NUK}+\textsf{NUB} (\textsf{S5}+\textsf{NUK}) (\textsf{K45}+\textsf{NUB}).
	\end{theorem}
	
	The last step is to define action models, both with and without revision, as in \cite{SimpAct} and \cite{TechMemo}. 
	
	\begin{definition}[Simplicial Action for Knowledge]
		A \defin{Simplicial Action for Knowledge} $A$ is a tuple $\langle N_A,V_A,L_A,S_A,Post\rangle$, where $Post:N_a\rightarrow 3^\mathfrak{P}$.\footnote{$N_A$ is a set of nodes, $V_A:N_A\rightarrow Ag$, $L_A:N_A\rightarrow 3^\mathfrak{P}$, and $S_A$ is a UCF subcomplex of $\mathfrak{M}(N_A,V_A,L_A)$, just as in a regular simplicial model.}
		
		Given a simplicial model for knowledge $\mathcal{M}=\langle N,V,L,S\rangle$, and a simplicial action $A=\langle N_A,V_A,L_A,S_A,Post\rangle$, we define the update model $\mathcal{M}[A]:=\langle N[A],V[A],L[A],S[A]\rangle$\footnote{$N[A]$ is a subset of $N\times N_A$, $V[A]:N[A]\rightarrow Ag$, $L[A]:N[A]\rightarrow 3^\mathfrak{P}$, and $S[A]$ is a UCF subcomplex of $\mathfrak{M}(N[A],V[A],L[A])$, just as in a regular simplicial model.} as follows:
		
		\begin{align*}
			N[A]&:=\{(x,y)\in N\times N_A|V(x)=V_A(y)\\
			&\wedge\forall P\in\mathfrak{P}((L(x)(P)=1\rightarrow L_A(y)(P)\neq 0)\wedge(L(x)(P)=0\rightarrow L_A(y)(P)\neq 1))\}
			\\ V[A]((x,y))&:=V(x)
			\\ L[A]((x,y))(P)&:=\begin{cases}
				& Post(y)(P)\;\text{if}\;Post(y)(P)\neq 2 \\
				& L_A(y)(P)\;\text{if}\;L(x)(P)=2 \\
				&L(x)\;\text{otherwise}
			\end{cases} 
			\\ \mathcal{F}(S[A])&:=\left \{F\in 2^{N[A]}~\middle |~\begin{matrix}
				& \pi_0(F)\in\mathcal{F}(S) \\
				& \wedge\pi_1(F)\in\mathcal{F}(S_A) \\
				& \wedge F\text{ is consistent w.r.t. }L[A] \\
				%& \wedge\forall x,x'\in\pi_0(F),y,y'\in\pi_1(F)(x=x'\leftrightarrow y=y')
			\end{matrix} \right \}
		\end{align*}
	\end{definition}
	
	Note that if $F\in 2^{N[A]}$, $\pi_0(F)=\{x~|~\exists y\in S_A((x,y))\in F)\}$, and similarly for $\pi_1$. Let us show that for any simplicial model for knowledge $\M$ and simplicial action $A$, $\M[A]$ satisfies the UCF property.\cite{SimpAct}
	
	Our definition of $N[A]$ is not the only one which could be motivated. In fact, all three of the definitions below are defensible.
	
	\begin{enumerate}
		\item $\forall P\in\mathfrak{P}((L(x)(P)=1\rightarrow L_A(y)(P)=1)\wedge(L(x)(P)=0\rightarrow L_A(y)(P)=0))$
		\item $\forall P\in\mathfrak{P}((L_A(y)(P)=1\rightarrow L(x)(P)=1)\wedge(L_A(y)(P)=0\rightarrow L_A(x)(P)=0))$
		\item $\forall P\in\mathfrak{P}((L(x)(P)=1\rightarrow L_A(y)(P)\neq 0)\wedge(L(x)(P)=0\rightarrow L_A(y)(P)\neq 1))$
	\end{enumerate}
	
	We will refer to these as modes 1-3. Indeed, mode 3 is the one introduced in our definition above. For our purposes in this paper, it will suffice.
	
	We now turn to modifying these notions for the belief setting:
	
	\begin{definition}[Simplicial Action for Belief]
		A \defin{Simplicial Action for Belief} $A$ is a tuple $\langle N_A,V_A,L_A,\{S_{a,A}\}_{a\in Ag},Post\rangle$, where $Post:N_a\rightarrow 3^\mathfrak{P}$.\footnote{$N_A$ is a set of nodes, $V_A:N_A\rightarrow Ag$, $L_A:N_A\rightarrow 3^\mathfrak{P}$, and each $S_{a,A}$ is a UCF subcomplex of $\mathfrak{M}(N_A,V_A,L_A)$, just as in a regular simplicial model.}
		
		Given a simplicial model for belief $\mathcal{M}=\langle N,V,L,\{S_{a}\}_{a\in Ag}\rangle$, And a simplicial action for belief $A=\langle N_A,V_A,L_A,\{S_{a,A}\}_{a\in Ag},Post\rangle$, we define the update model $\mathcal{M}[A]_B:=\langle N[A]_B,V[A]_B,L[A]_B,\{S_a[A]_B\}_{a\in Ag}\rangle$\footnote{$N[A]_B$ is a subset of $N\times N_A$, $V[A]_B:N[A]_B\rightarrow Ag$, $L[A]_B:N[A]_B\rightarrow 3^\mathfrak{P}$, and $S_a[A]_B$ is a UCF subcomplex of $\mathfrak{M}(N[A]_B,V[A]_B,L[A]_B)$, just as in a regular simplicial model.} as follows:
		
		\begin{align*}
			N[A]_B&:=\{(x,y)\in N\times N_A|V(x)=V_A(y)\\
			&\wedge\forall P\in\mathfrak{P}((L(x)(P)=1\rightarrow L_A(y)(P)\neq 0)\wedge(L(x)(P)=0\rightarrow L_A(y)(P)\neq 1))\}
			\\ V[A]_B((x,y))&:=V(x)
			\\ L[A]_B((x,y))(P)&:=\begin{cases}
				& Post(y)(P)\;\text{if}\;Post(y)(P)\neq 2 \\
				& L_A(y)(P)\;\text{if}\;L(x)(P)=2 \\
				&L(x)\;\text{otherwise}
			\end{cases} 
			\\ \mathcal{F}(S_a[A]_B)&:=\left \{F\in 2^{N[A]_B}~\middle |~\begin{matrix}
				& \pi_0(F)\in\mathcal{F}(S_a) \\
				& \wedge\pi_1(F)\in\mathcal{F}(S_{a,A}) \\
				& \wedge F\text{ is consistent w.r.t. }L[A]_B \\
				%& \wedge\forall x,x'\in\pi_0(F),y,y'\in\pi_1(F)(x=x'\leftrightarrow y=y')
			\end{matrix} \right \}
		\end{align*}
	\end{definition}

	We can also incorporate the notion of revision from \cite{SimpBelRev}, as discussed in \cite{SimpAct} and \cite{TechMemo}. The key idea behind revision, following \cite{Lewis2} and \cite{Stalnaker1981}, is that worlds are replaced by sets of worlds using a revision function. Since worlds are facets, we need to define a replacement function that maps facets in the input complex to sets of facets in the output of applying the action.

	\begin{definition}[Simplicial Action for Belief with Revision]
		Given a simplicial model for belief $\mathcal{M}=\langle N,V,L,\{S_{a}\}_{a\in Ag}\rangle$, and a simplicial action for belief $A=\langle N_A,V_A,L_A,\{S_{a,A}\}_{a\in Ag},Post\rangle$, we define the update model \defin{with revision} $\mathcal{M}[A]_{BR}:=\langle N[A]_{BR},V[A]_{BR},L[A]_{BR},\{S_a[A]_{BR}\}_{a\in Ag}\rangle$\footnote{$N[A]_B$ is a subset of $N\times N_A$, $V[A]_{BR}:N[A]_{BR}\rightarrow Ag$, $L[A]_{BR}:N[A]_{BR}\rightarrow 3^\mathfrak{P}$, and $S_a[A]_{BR}$ is a UCF subcomplex of $\mathfrak{M}(N[A]_{BR},V[A]_{BR},L[A]_{BR})$, just as in a regular simplicial model.} as follows:
		
		\begin{align*}
			N[A]_{BR}&:=\{(x,y)\in N\times N_A|V(x)=V_A(y)\\
			&\wedge\forall P\in\mathfrak{P}((L(x)(P)=1\rightarrow L_A(y)(P)\neq 0)\wedge(L(x)(P)=0\rightarrow L_A(y)(P)\neq 1))\}
			\\ V[A]_{BR}((x,y))&:=V(x)
			\\ L[A]_{BR}((x,y))(P)&:=\begin{cases}
				& Post(y)(P)\;\text{if}\;Post(y)(P)\neq 2 \\
				& L_A(y)(P)\;\text{if}\;L(x)(P)=2 \\
				&L(x)\;\text{otherwise}
			\end{cases} 
		\end{align*}
		
		Defining $S_a[A]_{BR}$ will take some more work. Fix $a\in Ag$ and let $X$ be a UCF facet in $2^{N\times N_A}$ such that $\pi_0(X)\in\mathcal{F}(S_a)$ and $\pi_1(X)\in\mathcal{F}(S_{a,A})$. Call the set of such $X$ $\mathcal{F}^a(N\times N_A)$. Let $X$ be a UCF facet in $2^{N[A]_{BR}}$ such that $X$ is consistent, and $\pi_1(X)\in\mathcal{F}(S_{a,A})$. Call the set of such $X$ $\mathcal{F}_a(2^{N[A]_{BR}})$. 
		
		Then we can define a replacement function as follows for $X\in\mathcal{F}^a(N\times N_A)$:
		
		$$R_a(X):=\{Y\in\mathcal{F}_a(2^{N[A]_{BR}})|\pi_a(Y)=\pi_a(X)$$
		$$\wedge(\forall Z\in\mathcal{F}_a(2^{N[A]_{BR}})(\pi_a(Z)=\pi_a(X)\rightarrow|Y\cap X|\geq|Z\cap X|))\}$$
		
		We say that $S_a[A]_{BR}$ is defined as the simplicial complex whose facets come from $\bigcup_{X\in\mathcal{F}^a(N\times N_A)}R_a(X)$.
	\end{definition}

	It is important to note that the application of an action to an input models a distributed protocol. Indeed, simplicial action models provably define a protocol in the sense of \cite{DCTCT}, \cite{SimpDEL}, and \cite{TopPersp}, as is shown in \cite{TechMemo} and \cite{SimpAct}.
	
	\bibliographystyle{plain}	
	\bibliography{Epistemology} 
	

	
\end{document}